\scanpage{49}

that these coefficients are rational, which implies that $w^{-1}$ is algebraic, and so is $w^{-1}u=v^{-1}$, which was to be proved.

NB In char. 0, \deleted{it is} the statement simplifies to: $C(X)$ equivalent to the existence of algebraic isomorphisms $H^{2n-i}(X)\to H^i(X)$,(as the preliminaray condition in b) is then automatically satisfied). Maybe with some extra care this can be proved too in arbitrary characteristics.

\uline{Corollary}\quad Assume $X$ and $X'$ satisfy condition $C_{\unsure{\ell|\chi}}$, and let $u:H^i(X)\mathrel{-}H^{i+2\unsure{\delta|d}}(X')$ \inserted{$(\unsure{\delta|d}\in\uline{Z})$} be an isomorphism which is algebraic. Then $u^{-1}$ is algebraic.

Indeed, the two spaces can be identified ``algebraically'' (both directions !) to their dual, so taht the transpose of $u$ can be viewed as an isomorphism $u':H^{i+2\unsure{\delta|d}}(X')\mathrel{-}H^i(X)$. Thus $u'u$ is an \inserted{algebraic} automorphism \inserted{$w$} of $H^i(X)$, and by the previous argument we see that $w^{-1}$ is algebraic, hence so is $u^{-1}=w^{-1}u'$.

As a consequence, we see that if $x\in H^i(X)$ is such that $u(x)$ is algebraic ($i$ being now assumed to be even), then so is $x$. The same result should hold in fact if $u$ is a monomorphism, the reason being that in this case there should exist a left-inverse which is algebraic; this exists indeed in a case like $H^{n-1}(X)\to H^{n-1}(Y)$ (where we take the left inverse $\Lambda_X\phi_*$). But to get \unsure{it} in general, it seems we need moreover the Hodge index relation. (The complete yoga then being that we have the category of motives which is semi-simple !). Without speeking of motives, and staying down on earth, it would be nice to explain in the notes that $C(X)$ together with the index relation $I(X\times X)$ implies that the ring of correspondance classes for $X$ is semi-simple, and how one deduces from this the existence \unsure{of|an} left and right inverses as looked for above.

\mspage{48}
